\section{Locked Branch-elimination Stress Audit}
\label{app:branch-elimination}

The countermodels were next subjected to every presently applicable locked constraint.  An alternative is rejected only if it violates an already-fixed physical ledger or fails a diagnostic whose candidate set and tolerance were fixed independently.  Residual ranking is not an exclusion rule.

The audit contains 103 candidate--test findings.  Eighty-four are physical rejections, one is an exact ledger rejection, one is a conditional physical rejection, nine tests are non-discriminating, and eight findings survive.  The
large rejection count is dominated by 61 electroweak coefficient triplets, 21 nonpublished single-monomial electromagnetic replacements, and the two nonminimal radial projectors $\Pi_{15}$ and $\Pi_{17}$.

\subsection{Compact Branches}

The fine-structure diagnostic rejects $\Pi_{15}$ and $\Pi_{17}$ by many orders of magnitude.  At fixed $\Pi_{13}$, however, all three terminal coefficients $\sigma_8=-1,0,+1$ survive.  The quasar observable
\begin{equation}
 \Delta v_{\rm high-low}
 =\operatorname{median}(\Delta v_{\rm est})_{\rm high\ L}
 -\operatorname{median}(\Delta v_{\rm est})_{\rm low\ L}
 \label{eq:pr35-quasar-nondiscrimination}
\end{equation}
contains no $\sigma_8$, $c_{\rm bc}$, or $\Pi$ dependence in its published form.  Each sign branch therefore makes the same negative-sign prediction. This successful PR-I stress test supports their common reduct but cannot select among its expansions. The Kerr residual bound \cite{Kerr_1963} is also unchanged when only $\sigma_8$ is varied at fixed $\Pi_{13}$, because the radial spectral gap is unchanged.  The magnetic law
\begin{equation}
 B_{\rm geo}^{\rm PR}
 =\frac{(\hbar/e)(Z_A/c_{\rm bc})}{A_{\rm proj}^{(\theta)}}
 \label{eq:pr35-magnetic-nondiscrimination}
\end{equation}
does inherit the small branch dependence of $c_{\rm bc}$.  Its present use is an inverse area constraint, however, and the realized source area is not fixed independently.  It consequently cannot exclude a terminal sign without a separate measurement or derivation of $A_{\rm proj}^{(\theta)}$.

\subsection{PR-II and PR-III Branches}

The vectorlike $t=0$ hypercharge family is excluded by the fixed nonzero electromagnetic charge ledger.  Complete generation replicas with $N_{\rm gen}\ne3$ are observationally nonphysical if sheet multiplicity is first identified with physical generation multiplicity, but every positive replication count continues to satisfy the local-anomaly and Witten gates. The missing step is therefore the identification theorem, not another anomaly calculation.

None of the 21 tested nonpublished single-monomial replacements for $D_1+D_2$ passes the locked fine-structure gate.  This eliminates that finite candidate family, but not arbitrary higher-order combinations.  Of the 64 tested electroweak triplets, only
\begin{equation}
 (a,b,k)=(2,-1,1),\quad(2,-1,2),\quad(2,-1,3)
\end{equation}
pass both locked mass gates.  The middle triplet is published and has the smallest $M_Z$ residual among these three integer cases, but the exact continuous central solution is $k=1.9567498823\ldots$, not two.  Choosing any coefficient on residual rank would be fitting, not a uniqueness proof.

The stress audit therefore achieves substantial physical elimination without producing a false singleton claim.  Appendices \ref{app:terminal-orientation-rigidity} and \ref{app:d3-coefficient-rigidity} now prove $\sigma_8=-1$ and the representation factor $k=2$ inside explicit boundary and unit-prefactor projected-adjoint classes.  The complete real-symmetric gauge-index coefficient classification at the fixed $D_3$ Lorentz/momentum kernel simultaneously proves that even after imposing the three stipulated structural gates every real $k$ remains within that matrix class.  Closure still requires deriving those classes, especially the pre-projection exclusion rule and the complete determinant-Hessian physical-$D_3$ operator map from the frozen construction, or obtaining independent observables whose equations genuinely separate the common-reduct branches.  The full matrix and reproducible classifications are provided with the source archive.
